\begin{pastexam}[label={px:ntu-108-bime-ds-06}]{臺大生機所 108 學年度 第 6 題}
（設計）

(15\%) How do you print all nodes of a given binary tree using inorder traversal without
recursion?
\end{pastexam}

\begin{pastsoln}
\textbf{想法：用自己的堆疊模擬遞迴}（資結書第~7.3.3~節虛擬碼~7.3）。 遞迴版進到一個節點時，要
先把左子樹整個做完才輪到它自己，所以「它自己」要先存起來。 迴圈版從根開始一路往左走，路上的
節點都 push；走到底之後 pop 出來的就是下一個要印的（它的左子樹已經做完），印出它，再轉到它的
右子樹，重複同樣的事。

\begin{pseudo}{中序走訪的迴圈版}
\Procedure{Inorder-Print}{$T$}
    \State 建立空堆疊 $S$
    \State $x \gets T.\mathit{root}$
    \While{$x \ne \NIL$ \OR $S$ 不是空的}
        \While{$x \ne \NIL$} \Comment{一路往左，路上的節點先存著}
            \State \Call{Push}{$S, x$}
            \State $x \gets x.\mathit{left}$
        \EndWhile
        \State $x \gets$ \Call{Pop}{$S$} \Comment{它的左子樹已經印完}
        \State 印出 $x.\mathit{key}$
        \State $x \gets x.\mathit{right}$
    \EndWhile
\EndProcedure
\end{pseudo}

\textbf{追蹤}，樹是 $A(B(D, E), C(F, \text{空}))$（根 $A$；$B$ 的孩子 $D$、$E$；$C$ 只有左孩子
$F$）。 每一列是一次印出，「推入」是印出前第 5--7 行 push 的節點：
\begin{center}
\begin{tabular}{ccl}
\toprule
印出 & 印出前推入 & 印出後的堆疊（底 $\to$ 頂） \\
\midrule
$D$ & $A, B, D$ & $A\;B$ \\
$B$ & & $A$ \\
$E$ & $E$ & $A$ \\
$A$ & & （空） \\
$F$ & $C, F$ & $C$ \\
$C$ & & （空） \\
\bottomrule
\end{tabular}
\end{center}
印出 $D\,B\,E\,A\,F\,C$，和遞迴版的中序一樣。 印完 $C$ 後 $x = C.\mathit{right} = \NIL$、堆疊
空，迴圈結束。

\textbf{正確性。} 每個節點恰好 push 一次、pop 一次。 一個節點 $x$ 被 pop 時，它左子樹的節點都
比它早 push、晚於它被 push 才 push 的節點都在它上面，已經先 pop 掉，所以左子樹已經全部印完；
接著印 $x$，再處理右子樹。 這正是「左、根、右」，和遞迴版每一步對應。

\textbf{複雜度。} 每個節點 push、pop 各一次，時間 $\Theta(n)$；堆疊最多同時放一條從根往下的
路徑，空間 $O(h)$，$h$ 是樹高。 如果連堆疊都不能用，可以用引線二元樹（資結書第~7.5~節）：沿著
引線找中序後繼，額外空間 $O(1)$。
\end{pastsoln}
